Mediante el criterio de Weyl, podemos demostrar facilmente que una sucesión es equidistribuida.
En este capítulo presentamos un método, inspirado en los ejercicios de Kuipers 
\cite[Ejercicios 2.21--2.22]{kuipers-1974}, para verificar dicha propiedad apoyándose en la fórmula de sumación de Euler. 

Esta fórmula nos permite expresar sumas de exponenciales en términos de integrales, las
cuales podemos acotar de forma precisa.


\begin{proposition}[Fórmula de sumación de Euler]
    Sea $N \in \mathbb{N}$ y sea $f:[1,N] \to \mathbb{C}$ una función tal que $f' \in C([1,N])$. 
    Entonces se cumple la siguiente igualdad:
    \begin{align}
        \label{prop:formula_euler}
        \sum_{n=1}^N f(n) = 
        \int_1^N f(x) \, dx + 
        \frac{f(1)+f(N)}{2} + 
        \int_1^N \left(\partfrac{x} - \frac{1}{2}\right) \, f'(x) \, dx
    \end{align}
\end{proposition}

\begin{proof}
    Dividimos el intervalo $[1,N]$ en subintervalos $[n,n+1]$ con $n=1,\ldots,N-1$. 
    
    Para cada intervalo $[n,n+1]$ se tiene que $\partfrac{x} = x - n$, 
    al integrar por partes,
    \begin{align*} 
        \int_n^{n+1} \left(\partfrac{x} - \frac{1}{2}\right) f'(x) \, dx 
        &= \left[\left(\partfrac{x} - \frac{1}{2}\right) f(x) \right]_n^{n+1} - \int_n^{n+1} f(x) \, dx \\
        &= \frac{1}{2} \, f(n+1) + \frac{1}{2} \, f(n) - \int_n^{n+1} f(x) \, dx
    \end{align*}

    Sumando sobre $n=1,2,\ldots,N-1$,
    \[
        \sum_{n=1}^{N-1} \int_n^{n+1} \left(\partfrac{x} - \frac{1}{2}\right) \, f'(x) \, dx =
        \sum_{n=1}^{N-1} \left(\frac{1}{2} \, f(n+1) + \frac{1}{2} \, f(n)\right) -
        \sum_{n=1}^{N-1} \int_n^{n+1} f(x) \, dx \\
    \]
    luego,
    \[
        \int_1^N \left(\partfrac{x} - \frac{1}{2}\right) \, f'(x) \, dx =
        \sum_{n=1}^N f(n) - \frac{1}{2} \, f(N) - \frac{1}{2} \, f(1) - 
        \int_1^N f(x) \, dx
    \]
    
    Reordenando los términos se obtiene la igualdad deseada.
\end{proof}

\begin{example}
    La sucesión $(a\log n)_{n\geq1}$ con $a \in \mathbb{R}$ no es equidistribuida módulo $1$.
\end{example}

\begin{resolve}
    Si $a = 0$, la secuencia es constante e igual a $0$, por lo que no es equidistribuida.
    
    Supongamos que $a \neq 0$. Verificaremos que no se cumple el criterio de Weyl para $k=1$. 
    Es decir,
    \[
        \lim_{N \to \infty} \frac{1}{N} \sum_{n=1}^N e^{2\pi i a \log n} \neq 0
    \]   
    
    Sea $f(x) = e^{2\pi i a \log x}$. Aplicando la fórmula (\ref{prop:formula_euler})
    \[
        \sum_{n=1}^N e^{2\pi i \, a \log n} =
        \int_1^N e^{2\pi i a \log x} \, dx +
        \frac{1}{2} \left(e^{2\pi i \, a \log N} + 1\right) +
        \int_1^N \left(\partfrac{x} - \frac{1}{2}\right) \frac{2\pi i a}{x} \, e^{2\pi i a \log x} \, dx 
    \]
    
    Calculamos la primera integral mediante el cambio de variable $u = \log x$,
    $x = e^u$ y $dx = e^u du$. Cuando $x = 1$, $u = 0$ y cuando $x = N$, $u = \log N$.
    \begin{align*}
        \int_1^N e^{2\pi i a \log x} \, dx 
        = \int_0^{\log N} e^{(2\pi i a + 1) u} \, du 
        &= \left[\frac{e^{(2\pi i a + 1) u}}{2\pi i a + 1}\right]_0^{\log N} \\ 
        &= \frac{e^{(2\pi i a + 1) \log N} - 1}{2\pi i a + 1} 
        = \frac{N e^{2\pi i a \log N} - 1}{2\pi i a + 1}
    \end{align*}
        
    Dividiendo por $N$, 
    \[
        \frac{N e^{2\pi i a \log N} - 1}{N(2\pi i a + 1)}
        = \frac{e^{2\pi i a \log N}}{2\pi i a + 1} - \frac{1}{N(2\pi i a + 1)}
    \]

    Cuando $N \to \infty$, el segundo término se anula, mientras que el primero 
    no converge, ya que oscila sobre la circunferencia unidad.

    El término central se puede acotar
    \[
        \left|\frac{1}{2} \left(e^{2\pi i a \log N} + 1\right)\right| \leq 1
    \]
    por lo que, tras dividir por $N$, támbien tiende a $0$.

    Para la segunda integral, usamos que $|\partfrac{x} - \frac{1}{2}| \leq \frac{1}{2}$ y que $|e^{2\pi i a \log x}| = 1$,
    \[
        \left|\int_1^N \left(\partfrac{x} - \frac{1}{2}\right) \dfrac{2\pi i a}{x} \, e^{2\pi i a \log x} \, dx\right|
        \leq \pi |a| \int_1^N \frac{1}{x} \, dx
        = \pi |a| \log N
    \]
    En consecuencia,
    \[
        \frac{1}{N} \left|\int_1^N \left(\partfrac{x} - \frac{1}{2}\right) \dfrac{2\pi i a}{x} \, e^{2\pi i a \log x} \, dx\right|
        \longrightarrow 0 \qquad \text{cuando} \ N \to \infty
    \]

    Como la primera integral no tiende a $0$, no se verifica el criterio de Weyl para $k=1$.
    Por tanto, la sucesión $(a\log n)$ no es equidistribuida módulo $1$.   
\end{resolve}


\begin{example}
    La sucesión $(an^\sigma)_{n\geq1}$ con $a \in \mathbb{R}$, $a \neq 0$ y 
    $\sigma \in (0,1)$ es equidistribuida módulo $1$.
\end{example}

\begin{resolve}    
    Para demostrarlo utilizaremos el criterio de Weyl. 

    Sea $k \in \mathbb{Z}\setminus\{0\}$, como $a \neq 0$, entonces $b = ak \neq 0$. 
    Basta comprobar que 
    \[
        \frac{1}{N} \sum_{n=1}^N e^{2\pi i b n^\sigma} \longrightarrow 0 
        \qquad \text{cuando} \ N \to \infty.
    \]

    Aplicando la fórmula (\ref{prop:formula_euler}) a la función $f(x) = e^{2\pi i b x^\sigma}$,
    \[
        \sum_{n=1}^N e^{2\pi i \, b n^\sigma}
        = \int_1^N e^{2\pi i b x^\sigma} \, dx 
        + \frac{1}{2} \left(e^{2\pi i \, b N^\sigma} + e^{2\pi i b}\right)
        + \int_1^N \left(\partfrac{x} - \frac{1}{2}\right) 2\pi i b \sigma x^{\sigma-1} \, e^{2\pi i b x^\sigma} \, dx 
    \]

    Calculamos la primera integral usando integración por partes,
    \begin{align*}
        \int_1^N e^{2\pi i b x^\sigma} \, dx 
        = \int_1^N \frac{\left(e^{2\pi i b x^\sigma}\right)'}{2\pi i b \sigma x^{\sigma-1}} \, dx
        &= \frac{1}{2 \pi i b \sigma} \int_1^N \left(e^{2\pi i b x^\sigma}\right)'x^{1-\sigma} \, dx \\
        &= \frac{1}{2 \pi i b \sigma} \left(\left[e^{2\pi i b x^\sigma}x^{1-\sigma}\right]_1^N + (1-\sigma) \int_1^N e^{2\pi i b x^\sigma} x^{-\sigma} \, dx \right)
    \end{align*}  
    
    Acotando cada uno de los términos anteriores,
    \begin{align*}
        \left| e^{2\pi i b N^\sigma}N^{1-\sigma} - e^{2\pi i b}\right|
        &\leq N^{1-\sigma} + 1 \\
        \left|\int_1^N e^{2\pi i b x^\sigma} x^{-\sigma} \, dx\right|
        &\leq \int_1^N x^{-\sigma} \, dx
        = \frac{N^{1-\sigma} - 1}{1-\sigma}
    \end{align*}
    Ambos crecen del orden de $N^{1-\sigma}$, por lo que la primera integral es$O(N^{1-\sigma})$.

    El término central se acota por
    \[
        \left|\frac{1}{2} \left(e^{2\pi i b N^\sigma} + e^{2\pi i b}\right)\right| \leq 1
    \]
    luego es $O(1)$.

    Para la segunda integral, usando que $|\partfrac{x} - \frac{1}{2}| \leq \frac{1}{2}$,
    y que $|e^{2\pi i b x^\sigma}| = 1$, se tiene
    \[
        \left|\int_1^N \left(\partfrac{x} - \frac{1}{2}\right) 2\pi i b \sigma x^{\sigma-1} \, e^{2\pi i b x^\sigma} \, dx\right|
        \leq \pi |b| \int_1^N \sigma x^{\sigma-1} \, dx  
        = \pi |b| (N^\sigma - 1)
    \]
    por lo que esta integral es $O(N^\sigma)$.
    
    Reuniendo las estimaciones anteriores, y teniendo en cuenta que 
    $\sigma \in (0,1)$, se obtiene
    \[
        \sum_{n=1}^N e^{2\pi i \, b n^\sigma} = O(N^{1-\sigma}) + O(1) + O(N^\sigma) = O(N^{\max\{\sigma, 1-\sigma\}}) = o(N)
    \]
    De este modo,
    \[
        \frac{1}{N} \sum_{n=1}^N e^{2\pi i b n^\sigma} \longrightarrow 0 
        \qquad \text{cuando} \ N \to \infty,
    \]
    y por tanto la sucesión $(an^\sigma)$ es equidistribuida módulo $1$.
\end{resolve}



La idea del ejemplo anterior se puede extender a funciones más generales que cumplan 
ciertas condiciones de crecimiento. El siguiente resultado, el teorema de Fejér, 
nos dice con precisión cuáles son estas condiciones.

    Mediante la fórmula de sumación de Euler (\ref{prop:SumacionEuler}), sea $f(x) = e^{2\pi i k g(x)}$ y sea $k \in \mathbb{Z} \setminus \{0\}$.
    Veamos que condiciones debe cumplir $g$ para que se cumpla el criterio de Weyl.
    \[
        \sum_{n=1}^N e^{2\pi i k g(n)}
        = \int_1^N e^{2\pi i k g(x)} \, dx 
        + \frac{1}{2} \left(e^{2\pi i k g(N)} + e^{2\pi i k g(1)}\right)
        + \int_1^N \left(\partfrac{x} - \frac{1}{2}\right) 2\pi i k g'(x) \, e^{2\pi i k g(x)} \, dx
    \]

    La primera integral la calculamos usando integración por partes.
    \begin{align*}
        \int_1^N e^{2\pi i k g(x)} \, dx 
        = \int_1^N \frac{\left(e^{2\pi i k g(x)}\right)'}{2\pi i k g'(x)} \, dx
        &= \frac{1}{2 \pi i k} \int_1^N \frac{\left(e^{2\pi i k g(x)}\right)'}{g'(x)} \, dx \\
        &= \frac{1}{2 \pi i k} \left(\left[\frac{e^{2\pi i k g(x)}}{    g'(x)}\right]_1^N + \int_1^N e^{2\pi i k g(x)} \left(\frac{1}{g'(x)}\right)' \, dx \right)
    \end{align*}
    Vamos a ver cada termino por separado.
    \begin{align*}
        \left|\frac{e^{2\pi i k g(N)}}{g'(N)} - \frac{e^{2\pi i k g(1)}}{g'(1)}\right|
        &\leq \frac{1}{|g'(N)|} + \frac{1}{|g'(1)|} \\
        \left|\int_1^N e^{2\pi i k g(x)}\left(\frac{1}{g'(x)}\right)' \, dx\right|
        \leq \int_1^N \left|\left(\frac{1}{g'(x)}\right)' \right| \, dx 
        &= \left|\frac{1}{g'(N)} - \frac{1}{g'(1)}\right|
        \leq \frac{1}{|g'(N)|} + \frac{1}{|g'(1)|}
    \end{align*}
    Por la monotonia de $g'$ podemos quitar el valor absoluto de la integral
    Por lo tanto, sea $C_1$ una constante que amontone las cotas todo lo que no depende de $N$.
    \[
        \left|\frac{1}{N} \int_{1}^N e^{2\pi i k g(n)}\right| \leq \frac{C_1}{N \left|g'(x)\right|}
    \]
    Necesitamos que se cumpla la condicion $\lim_{x \to \infty} {x \left|g'(x)\right|} = \infty$

    El termino central se puede acotar por $1$, no necesitamos nada especial para él.

    La última integral se puede acotar usando que $|\partfrac{x} - \frac{1}{2}| \leq \frac{1}{2}$, $|e^{2\pi i k g(x)}| = 1$ y por lo tanto
    \[
        \left|\int_1^N \left(\partfrac{x} - \frac{1}{2}\right) 2\pi i k g'(x) \, e^{2\pi i k g(x)} \, dx\right|
        \leq \int_1^N \frac{1}{2} \, 2\pi |k| |g'(x)| \, dx
        = \pi |k| (g(N) - g(1))
    \]
    Aglutinando todo en otra constante $C_2$ que amontone las cotas todo lo que no depende de $N$,
    \[
        \left|\frac{1}{N} \int_1^N \left(\partfrac{x} - \frac{1}{2}\right) 2\pi i k g'(x) \, e^{2\pi i k g(x)} \, dx\right| \leq C_2 \left|\frac{g(N)}{N}\right|
    \]
    Necesitamos que se cumpla la condición $\lim_{x \to \infty} \frac{g(x)}{x} = 0$.

\begin{example}
    Las sucesiones $(an^\sigma\log^\tau n)_{n\geq2}$ con $a \neq 0$, 
    $\sigma \in (0,1)$ y $\tau \in \mathbb{R}$ son equidistribuidas módulo $1$.
\end{example}

\begin{resolve}
    Sea $x_n = an^\sigma\log^\tau n$ para cada $n \geq 2$. Definimos $g(x) = a(x+1)^\sigma\log^\tau (x+1)$
    en~$[1,\infty)$, entonces $g(n) = x_{n+1}$ para cada $n \geq 1$, y su derivada es 
    \[
        g'(x) = a (x+1)^{\sigma-1} \log^{\tau-1}(x+1) \left(\sigma \log(x+1) + \tau\right)
    \]
    esta es continua en $[1,\infty)$ por ser composición de continuas y como $\sigma \in (0,1)$, 
    \[
        \lim_{x \to \infty} g'(x) = 0
        \qquad \text{y} \qquad
        \lim_{x \to \infty} x |g'(x)| = \infty
    \]

    Monotonia, estudiar como hacerlo.

    Por el teorema \ref{teo:fejer_euler}, $(g(n))$ es equidistribuida módulo $1$ y 
    por lo tanto, $(x_n)$ también lo es.
\end{resolve}

\begin{example}
    Las sucesiones $(a\log^\tau n)_{n\geq1}$ con $a \neq 0$ y $\tau > 1$ son equidistribuidas módulo $1$.
\end{example}

\begin{resolve}
    Sea $g(x) = a\log^\tau (x)$ entonces su derivada es
    \[
        g'(x) = \frac{a \tau \log^{\tau-1}(x)}{x}
    \]
    como $\tau > 1$, es continua en $[1,\infty)$ por ser composición de continuas. Además, 
    \[
        \lim_{x \to \infty} g'(x) = 0
        \qquad \text{y} \qquad
        \lim_{x \to \infty} x |g'(x)| = \infty
    \]

    Para ver la monotonía de $g'(x)$, calculamos su derivada,
    \[
        g''(x) = a \tau \frac{\log^{\tau-2}(x)}{x^2} \left((\tau - 1) - \log(x)\right)
    \]
    Obtenemos que $g''(x)=0$ en $x = e^{\tau - 1}$, por lo que, para $x > e^{\tau - 1}$ tendremos que
    $g''(x)$ es de signo constante. Luego $g'(x)$ es monótona a partir de $x_0 = e^{\tau - 1}$.

    Por el teorema \ref{teo:fejer_euler}, la sucesión $(a\log^\tau (n))$ es equidistribuida módulo $1$.
\end{resolve}

\begin{example}
    Las sucesiones $(an\log^\tau n)_{n\geq2}$ con $a \neq 0$ y $\tau < 0$ son equidistribuidas módulo $1$.
\end{example}

\begin{resolve}
    Sea $x_n = an\log^\tau n$ para cada $n \geq 2$. Definimos $g(x) = a(x+1)\log^\tau (x+1)$
    en~$[1,\infty)$, entonces $g(n) = x_{n+1}$ para cada $n \geq 1$, y su derivada es 
    \[
        g'(x) = a \log^{\tau}(x+1) + a \tau \log^{\tau-1}(x+1)
    \]
    que es continua en $[1,\infty)$ por ser composición de continuas y como $\tau < 0$, 
    \[
        \lim_{x \to \infty} g'(x) = 0
        \qquad \text{y} \qquad
        \lim_{x \to \infty} x |g'(x)| = \infty
    \]

    Para comprobar la monotonía de $g'(x)$, calculamos su derivada,
    \[
        g''(x) = a \tau \frac{\log^{\tau-2}(x+1)}{x+1} \left(\log(x+1) + \tau - 1\right)
    \]
    Como $g''(x)=0$ en $x = e^{1-\tau} - 1$, entonces para $x \geq e^{1-\tau} - 1$ 
    se tiene que $g''(x)$ es de signo constante. En consecuencia, $g'(x)$ es monótona a 
    partir de $x_0 = e^{1-\tau} - 1$.

    Por el teorema \ref{teo:fejer_euler}, $(g(n))$ es equidistribuida módulo $1$ y 
    por lo tanto, $(x_n)$ también lo es.
\end{resolve}